% 海王星（综述）
% license CCBYSA3
% type Wiki

本文根据 CC-BY-SA 协议转载翻译自维基百科\href{https://en.wikipedia.org/wiki/Neptune}{相关文章}。

\begin{figure}[ht]
\centering
\includegraphics[width=8cm]{./figures/bfb6ab5163cda4f5.png}
\caption{1989年8月由旅行者2号拍摄的真实色彩海王星图像；\(^\text{[1]}\) 中央为大黑斑\(^\text{[a]}\)。} \label{fig_Neptun_1}
\end{figure}
海王星是绕太阳运行的第八颗也是已知最远的行星。按直径计算，它是太阳系中第四大行星，按质量计算为第三大行星，并且是密度最高的巨行星。它的质量是地球的17倍。与其邻近的冰巨行星天王星相比，海王星略小，但质量更大、密度更高。由于主要由气体和液体组成\(^\text{[21]}\)，它没有明确定义的固体表面。海王星每164.8年绕太阳运行一周，轨道距离为30.1个天文单位（45亿千米；28亿英里）。它以罗马海神命名，并拥有天文学符号♆，代表海神的三叉戟\(^\text{[e]}\)。

海王星肉眼不可见，也是太阳系中唯一一颗最初不是通过直接经验观测而发现的行星。相反，天王星轨道的异常变化使亚历克西·布瓦尔提出假说，认为其轨道受未知行星的引力扰动。在布瓦尔去世后，约翰·库奇·亚当斯和乌尔班·勒维耶分别依据布瓦尔的观测独立地数学预测了海王星的位置。随后在1846年9月23日\(^\text{[2]}\)，约翰·戈特弗里德·加勒用望远镜直接观测到了海王星，位置与勒维耶预测的位置相差不到一度。其最大卫星海卫一不久后被发现，但海王星剩余的其他卫星直到20世纪才通过望远镜观测到。

海王星与地球的距离使其视直径很小，而其距离太阳遥远则使其极为暗淡，给基于地面的望远镜研究带来挑战。直到哈勃太空望远镜以及配备自适应光学的大型地面望远镜出现，才使得详细观测成为可能。1989年8月25日飞掠海王星的旅行者2号仍然是唯一访问过这颗行星的航天器\(^\text{[22][23]}\)。与气体巨行星（木星和土星）类似，海王星的大气主要由氢和氦组成，并含有少量的碳氢化合物以及可能存在的氮，但其含有更高比例的冰，如水、氨和甲烷。与天王星类似，其内部主要由冰和岩石构成\(^\text{[24]}\)，这两颗行星通常被视为“冰巨行星”以示区分\(^\text{[25]}\)。除了瑞利散射外，外层区域中微量的甲烷使海王星呈现出淡蓝色\(^\text{[26][27]}\)。

与天王星强烈的季节性大气不同，后者可能长时间几乎没有可见特征，海王星的大气具有活跃且持续可见的天气形态。在旅行者2号1989年飞掠期间，该行星的南半球存在一个可与木星大红斑相比的大黑斑。2018年，新的主要暗斑及更小的暗斑被识别并研究\(^\text{[28]}\)。这些天气形态由太阳系所有行星中最强的持续风驱动，风速可高达 2{,}100 km/h（580 m/s；1{,}300 mph）\(^\text{[29]}\)。由于其与太阳的巨大距离，海王星外层大气是太阳系中最寒冷的地方之一，其云顶温度接近 55 K（−218 °C；−361 °F）。行星中心温度约为 5{,}400 K（5{,}100 °C；9{,}300 °F）\(^\text{[30][31]}\)。海王星具有微弱且破碎的环系统（称为“弧”），于 1984 年被发现并由旅行者2号确认\(^\text{[32]}\)。
\subsection{历史}
\subsubsection{发现}
\begin{figure}[ht]
\centering
\includegraphics[width=6cm]{./figures/b8d92d1f5c26ae9a.png}
\caption{伽利略·伽利莱可能是第一位观测到海王星的人} \label{fig_Neptun_2}
\end{figure}
\begin{figure}[ht]
\centering
\includegraphics[width=6cm]{./figures/a7ed05f7bd7b8f95.png}
\caption{乌尔班·勒维耶大体上成功预测了海王星的位置} \label{fig_Neptun_3}
\end{figure}
\begin{figure}[ht]
\centering
\includegraphics[width=6cm]{./figures/bb642f227f35d3c1.png}
\caption{约翰·戈特弗里德·加勒受勒维耶之请求寻找海王星} \label{fig_Neptun_4}
\end{figure}
\begin{figure}[ht]
\centering
\includegraphics[width=6cm]{./figures/6e8babe86270c7ba.png}
\caption{约翰·库奇·亚当斯独立计算了海王星的位置} \label{fig_Neptun_5}
\end{figure}
一些已知最早的望远镜观测记录是伽利略于1612年12月28日和1613年1月27日（新历）绘制的图，它们包含的定位点与现今已知的海王星在这些日期的位置相吻合。两次观测中，伽利略似乎都把海王星误认为一颗恒星，因为它在夜空中出现在靠近木星的位置——即形成合相\(^\text{[33]}\)。因此，他不被视为海王星的发现者。在1612年12月的首次观测中，海王星在天空中几乎保持静止，因为它就在那一天开始逆行。这种视逆行运动是当地球轨道超过外行星时产生的。由于海王星只是刚刚开始其每年的逆行周期，这颗行星的运动非常微弱，以至于伽利略的小型望远镜无法检测到\(^\text{[34]}\)。2009年的一项研究表明，伽利略至少意识到他所观测到的这颗“恒星”相对于固定恒星发生了移动\(^\text{[35]}\)。

1821年，亚历克西·布瓦尔发表了关于天王星轨道的天文表\(^\text{[36]}\)。随后观测揭示出与天文表相比存在显著偏差，这促使布瓦尔假设有未知天体通过引力相互作用扰动了天王星的轨道\(^\text{[37]}\)。1843年，约翰·库奇·亚当斯开始利用他所拥有的数据研究天王星的轨道。他向皇家天文学家乔治·艾里爵士请求额外数据，艾里于1844年2月提供了这些数据。亚当斯在1845—1846年继续工作，并提出了多个关于天王星之外未发现行星位置的不同估算\(^\text{[38][39]}\)。

与亚当斯独立地，乌尔班·勒维耶在1845—1846年自行得出了指向未发现行星的计算结果，但并未激起其同胞的热情。1846年6月，在看到勒维耶首次发表的可疑未发现行星经度估算及其与亚当斯估算的相似性后，艾里说服詹姆斯·查利斯进行搜索。查利斯在整个8月和9月徒劳地搜索天空\(^\text{[37][40]}\)。事实上，查利斯比后来真正发现该行星的约翰·戈特弗里德·加勒早一年观测到了海王星，并在1845年8月4日和12日两次观测到。然而，他过时的星图和糟糕的观测技术使他在后续分析之前未能识别这些观测结果。查利斯充满悔意，但将自己的疏忽归咎于星图以及同时进行彗星观测工作的干扰\(^\text{[41][37][42]}\)。

与此同时，勒维耶寄出了一封信，并敦促柏林天文台天文学家加勒使用天文台的折射望远镜进行搜索。海因里希·德阿雷斯特（Heinrich d'Arrest），一位在天文台学习的学生，向加勒建议，他们可以将雷维耶预测位置区域的最新绘制星图与当前星空进行比较，以寻找行星特有的位移，而非恒星那样的固定位置。1846年9月23日傍晚，也就是加勒收到信件的当天，他就在宝瓶座ι星东北方向、距离勒维耶所预测的“摩羯座δ星以东五度”位置1°处发现了海王星\(^\text{[43][44]}\)，该位置距离亚当斯的预测约12°，并根据现代国际天文学联合会（IAU）星座边界处于宝瓶座与摩羯座的交界处。
\begin{figure}[ht]
\centering
\includegraphics[width=6cm]{./figures/80ca30e615e65bcf.png}
\caption{加勒用于发现海王星的9英寸折射望远镜} \label{fig_Neptun_6}
\end{figure}
在发现之后，法国人与英国人之间围绕谁应当获得这一发现的荣誉产生了民族主义式的竞争。最终，国际共识形成，认为勒维耶和亚当斯应当共同享有荣誉\(^\text{[45]}\)。自1966年以来，丹尼斯·罗林斯（Dennis Rawlins）一直质疑亚当斯共同发现主张的可信性，而随着1998年“海王星文献”（历史文件）归还格林尼治皇家天文台，该问题又被历史学家重新评估\(^\text{[46][47]}\)。
\subsubsection{命名}
在其发现后不久，海王星被简单地称为“天王星外侧的行星”或“勒维耶的行星”。第一个命名建议来自加勒，他提议命名为“雅努斯”。在英国，查利斯提出了“欧绪安努斯”这一名称\(^\text{[48]}\)。

勒维耶宣称自己拥有为其发现命名的权利，迅速为这颗新行星提出了“海王星”这一名称，尽管他虚假声称该名称已获得法国经度局的正式批准\(^\text{[49]}\)。在十月，他试图以自己的名字“勒维耶”命名该行星，并得到了天文台台长弗朗索瓦·阿拉戈的忠实支持。这一建议在法国之外遭到了强烈抵制\(^\text{[50]}\)。法国年鉴很快重新启用了“赫歇尔”一名来指代天王星，以纪念该行星的发现者威廉·赫歇尔爵士，并用“勒维耶”来指代新行星\(^\text{[51]}\)。

斯特鲁维于1846年12月29日在圣彼得堡科学院赞成使用“海王星”这一名称\(^\text{[52]}\)，这一命名是基于通过望远镜观察到的行星颜色\(^\text{[53]}\)。很快，“海王星”成为国际公认的名称。在罗马神话中，海王星是海神，与希腊神话中的波塞冬相对应。对神话名称的需求似乎与其他行星的命名法相一致，后者均以希腊和罗马神话中的神祇命名\(^\text{[f][54]}\)。

今天，大多数语言都使用“海王星”名称的某种变体。在中文、越南语、日语和韩语中，该行星名称被翻译为“海王星”\(^\text{[55][56]}\)。在蒙古语中，海王星被称为“Dalain van”（Далайн ван），反映其同名神祇作为海之统治者的角色。在现代希腊语中，该行星被称为“Poseidon”（Ποσειδώνας, Poseidonas），即希腊对应神波塞冬\(^\text{[57]}\)。在希伯来语中，“Rahab”（רהב）来自《诗篇》中提到的圣经海怪，在2009年由希伯来语言学院管理的投票中被选为该行星的官方名称，尽管现有的拉丁语形式“Neptun”（נפטון）更常使用\(^\text{[58][59]}\)。在毛利语中，该行星被称为“Tangaroa”，取自毛利海神之名\(^\text{[60]}\)。在纳瓦特尔语中，该行星被称为“Tlāloccītlalli”，取自雨神特拉洛克之名\(^\text{[60]}\)。在泰语中，海王星被称为西化名称“Dao Nepchun/Nepjun”（ดาวเนปจูน），但也被称为“Dao Ket”（ดาวเกตุ，意为“计都之星”），取自印度占星术中扮演角色的下降交点计都（केतु）。在马来语中，“Waruna”这一名称自20世纪70年代起有使用\(^\text{[61]}\)，取自印度海神之名，但最终被拉丁语形式“Neptun”（马来西亚用法\(^\text{[62]}\)）或“Neptunus”（印度尼西亚用法\(^\text{[63]}\)）所取代。

常用的形容词形式为“Neptunian”。一度出现过“Poseidean”（/pəˈsaɪdiən/）这一临时形式，取自“Poseidon”，尽管“Poseidon”的常用形容词形式是“Poseidonian”（/ˌpɒsaɪˈdoʊniən/）\(^\text{[5][64]}\)。
\subsubsection{地位}
从1846年被发现起直到1930年冥王星被发现之前，海王星是已知的最远行星。冥王星被发现后被视为行星，因此海王星成为第二远的已知行星，唯有在1979年至1999年这20年间，由于冥王星的椭圆轨道使其比海王星更靠近太阳，使海王星在此期间成为距离太阳第九颗行星\(^\text{[65][66]}\)。对冥王星质量的估算越来越准确，从地球的十倍降至远低于月球，并且1992年柯伊伯带的发现，使许多天文学家开始讨论冥王星应被视为行星还是柯伊伯带的一部分\(^\text{[67][68][69]}\)。2006年，国际天文学联盟首次定义了“行星”一词，将冥王星重新分类为“矮行星”，使海王星再次成为太阳系中已知最外侧的行星\(^\text{[70]}\)。
\subsection{物理特征}
\begin{figure}[ht]
\centering
\includegraphics[width=8cm]{./figures/dbcd02bcd550af2d.png}
\caption{海王星与地球的尺寸比较} \label{fig_Neptun_7}
\end{figure}
海王星的质量为 \(1.024\times10^{26}\,\text{kg}\)\(^\text{[8]}\)，介于地球和更大的气体巨行星之间：其质量是地球的17.15倍，但仅为木星的1/19\(^\text{[g]}\)。在1巴压力处，其重力为 \(11.27\,\text{m/s}^2\)，是地球表面重力的1.15倍\(^\text{[8]}\)，仅次于木星\(^\text{[71]}\)。海王星的赤道半径为 \(24{,}764\,\text{km}\)\(^\text{[11]}\)，几乎是地球的四倍。海王星与天王星一样，是冰巨行星（巨行星的一个子类），因为它们比木星和土星更小，并且含有更高浓度的挥发物\(^\text{[72]}\)。在寻找系外行星的研究中，海王星常被用作转喻：那些发现的质量相近的天体通常被称为“海王星类”\(^\text{[73]}\)，就像科学家将各种系外天体称为“类木星”一样。
\subsubsection{内部结构}
海王星的内部结构类似于天王星。其大气占其质量的约5\%到10\%，并向内部延伸至距核心大约10\%到20\%的位置。大气层中的压力约达到 \(10\,\text{GPa}\)，即约 \(10^5\) 个大气压。大气层较低区域中含有更高浓度的甲烷、氨和水\(^\text{[30]}\)。
\begin{figure}[ht]
\centering
\includegraphics[width=8cm]{./figures/5b0e22d5659369d9.png}
\caption{以伪彩色呈现的海王星内部物理组成及其周围结构示意图} \label{fig_Neptun_8}
\end{figure}
地幔的质量相当于10到15个地球质量，且富含水、氨和甲烷\(^\text{[2]}\)。按照行星科学的惯例，这种混合物被称为“冰质”，尽管它实际上是一种高温、致密的超临界流体。这种具有高电导率的流体有时被称为“水—氨海”\(^\text{[74]}\)。地幔可能由一层离子水组成，其中水分子分解为氢离子和氧离子的“汤”；在更深处则存在超离子水，其中氧形成晶格，但氢离子在氧晶格内自由移动\(^\text{[75]}\)。在7{,}000公里深处，条件可能使甲烷分解成钻石晶体，并像冰雹一样向下坠落\(^\text{[76][77][78]}\)。科学家认为，这种“钻石雨”也出现在木星、土星和天王星上\(^\text{[79][77]}\)。劳伦斯利弗莫尔国家实验室的超高压实验表明，地幔顶部可能是一片由液态碳构成的“海洋”，其中漂浮着固态“钻石”\(^\text{[80][81][82]}\)。

海王星的核心很可能由铁、镍和硅酸盐组成，内部模型显示其质量约为地球的1.2倍\(^\text{[24]}\)。中心压力为7兆巴（700 GPa），约为地球中心压力的两倍，温度可能为5{,}400 K（5{,}100 °C；9{,}300 °F）\(^\text{[30][31]}\)。
\subsubsection{大气}
\begin{figure}[ht]
\centering
\includegraphics[width=8cm]{./figures/247ec18a4a8a7802.png}
\caption{1994年至2020年的海王星云层变化\(^\text{[83]}\)。该伪彩色图像基于哈勃太空望远镜WFPC2与WFC3仪器的数据。} \label{fig_Neptun_9}
\end{figure}
\begin{figure}[ht]
\centering
\includegraphics[width=8cm]{./figures/8072a9d15d9dd52f.png}
\caption{高空云带在海王星较低层的云面上投下阴影。为使云层更清晰，图像颜色被夸张处理。} \label{fig_Neptun_10}
\end{figure}
在高空，海王星大气由80\%的氢和19\%的氦组成\(^\text{[30]}\)，并含有微量甲烷。甲烷在光谱600\,\text{nm}以上的波长（红光和红外区）具有显著的吸收带。与天王星一样，大气甲烷对红光的吸收是造成海王星微弱蓝色调的原因之一\(^\text{[84]}\)，但由于天王星大气中浓集的薄雾，海王星的蓝色更为明显\(^\text{[85][86]}\)。

海王星大气分为两个主要区域：随高度上升温度下降的对流层，以及随高度上升温度升高的平流层。二者之间的边界，即对流层顶，位于0.1\,\text{bar}（10\,\text{kPa}）的压力处\(^\text{[25]}\)。随后平流层在低于 \(10^{-5}\) 至 \(10^{-4}\,\text{bar}\)（1至10\,\text{Pa}）的压力下过渡至热层\(^\text{[25]}\)。热层逐渐过渡到外逸层\(^\text{[87]}\)。

模型显示，海王星的对流层中存在不同成分的带状云层，其成分随高度改变\(^\text{[83]}\)。上层云位于低于1\,\text{bar}的压力处，在该处温度适合于甲烷凝结。介于1至5\,\text{bar}（100至500\,\text{kPa}）之间，认为会形成氨和硫化氢云层。在超过5\,\text{bar}压力之上，云层可能由氨、硫氨、硫化氢和水组成。约50\,\text{bar}（5.0\,\text{MPa}）压力处应存在水冰云层，在此温度达到273\,\text{K}（0\,^\circ\text{C}；32\,^\circ\text{F}）。在更下方可能存在氨和硫化氢云层\(^\text{[88]}\)。

已观测到海王星的高空云层在下方不透明云层上投下阴影。存在环绕行星、纬度固定的高空云带。这些环状带宽度为50–150\,\text{km}，位于云面上方约50–110\,\text{km}处\(^\text{[89]}\)。这些高度属于发生天气活动的对流层；天气并不发生在更高的平流层或热层。2023年8月，海王星高空云层消失，引发一项历时三十年的研究，该研究结合了哈勃太空望远镜和地面望远镜的观测。研究表明，海王星的高空云活动与太阳周期相关，而非与行星季节相关\(^\text{[83][90][91]}\)。

海王星光谱显示，其较低平流层由于甲烷紫外光解产物（如乙烷和乙炔）的凝结而呈薄雾状\(^\text{[25][30]}\)。平流层中存在微量一氧化碳和氰化氢\(^\text{[25][92]}\)。由于更高浓度的烃类，海王星平流层比天王星更温暖\(^\text{[25]}\)。

目前原因仍不清楚，行星热层具有约750\,K（$477\,^\circ\text{C}$；$890\,^\circ\text{F}$）的异常高温\(^\text{[93][94]}\)。距离太阳过远，无法由紫外辐射产生如此热量。一种候选加热机制是大气与行星磁场中离子的相互作用。其他候选机制包括从内部传播并在大气中耗散的重力波。热层中含有微量二氧化碳和水，这些物质可能来源于如流星体和尘埃等外部来源\(^\text{[88][92]}\)。
\subsubsection{颜色}
海王星的大气在可见光光谱中呈微弱蓝色，仅比天王星大气的蓝色稍微饱和一些。早期两颗行星的图像大大夸张了海王星的色彩对比“以更好地呈现云层、带状结构和风”，使其相比天王星的灰白色显得深蓝。由于两颗行星使用了不同成像系统进行拍摄，很难直接比较后期合成图像。随着时间推移，这一问题得以重新评估，并在2023年底进行了最全面的颜色归一化\(^\text{[86][95]}\)。











